\documentclass[12pt]{article}
\usepackage[a4paper, margin=1in]{geometry}
\usepackage{amsmath, amssymb}
\usepackage{enumitem}
\usepackage{hyperref}

\begin{document}

\section*{\textbf{CSE400: Fundamentals of Probability in Computing}}

\subsection*{\textbf{Lecture Scribe 9: Uniform, Exponential, Laplace and Gamma Random Variables}}

\textbf{Instructor:} Dhaval Patel, PhD\\
\textbf{Date:} February 2, 2026

\hrule
\vspace{1em}

\section*{\textbf{1. Outline}}

This lecture covers the following topics:

\begin{itemize}
    \item Types of Continuous Random Variables
    \item Uniform Random Variable: Example
    \item Exponential Random Variable: Example
    \item Laplace Random Variable: Example
    \item Gamma Random Variable
    \begin{itemize}
        \item Graph and Special Cases
        \item Example
        \item Homework Problem
    \end{itemize}
    \item Problem Solving
    \item In-Class Activity: Gaussian Density Estimation
\end{itemize}

\hrule
\vspace{1em}

\section*{\textbf{2. Types of Continuous Random Variables}}

\subsection*{\textbf{2.1 Uniform Random Variable}}

For a uniform random variable on the interval $([a,b])$, where $(a<x<b)$:

\subsubsection*{\textbf{PDF}}

\[
f_X(x)=
\begin{cases}
\frac{1}{b-a}, & a<x<b \\
0, & \text{elsewhere}
\end{cases}
\]

\subsubsection*{\textbf{CDF}}

\[
F_X(x)=
\begin{cases}
0, & x<a \\
\frac{x-a}{b-a}, & a<x<b \\
1, & x>b
\end{cases}
\]

\hrule
\vspace{1em}

\subsection*{\textbf{2.2 Example \#1 (Uniform RV)}}

\subsubsection*{\textbf{Problem}}

The phase of a sinusoid, $(\Theta)$, is uniformly distributed over $([0,2\pi))$ so that its PDF is:

\[
f_\Theta(\theta)=
\begin{cases}
\frac{1}{2\pi}, & 0<\theta<2\pi \\
0, & \text{otherwise}
\end{cases}
\]

Find:

\begin{itemize}
    \item (a) $(\Pr(\Theta > 3\pi))$
    \item (b) $(\Pr(\Theta < \pi \mid \Theta > 3\pi))$
    \item (c) $(\Pr(\cos(\Theta) < \frac{1}{2}))$
\end{itemize}

\hrule
\vspace{1em}

\subsubsection*{\textbf{Solution: Example 1 (Uniform RV)}}

Given:

\[
f_\Theta(\theta)=
\begin{cases}
\frac{1}{2\pi}, & 0<\theta<2\pi \\
0, & \text{otherwise}
\end{cases}
\]

For a uniform RV on $([0,2\pi))$:

\[
\Pr(a<\Theta<b)=\frac{b-a}{2\pi}
\]

\subsubsection*{\textbf{(a) Find $(\Pr(\Theta > 3\pi))$}}

\[
\Pr(\Theta > 3\pi)=\Pr(\Theta>\pi)=\frac{2\pi-\pi}{2\pi}=\frac{1}{2}
\]

\subsubsection*{\textbf{(b) Find $(\Pr(\Theta < \pi \mid \Theta > 3\pi))$}}

\[
\Pr(A\mid B)=\frac{\Pr(A\cap B)}{\Pr(B)}
\]

\subsubsection*{\textbf{(c) Find $(\Pr(\cos(\Theta)<\frac{1}{2}))$}}

\[
\cos\Theta=\frac{1}{2}\Rightarrow \Theta=\frac{\pi}{3}
\]

\[
\cos\Theta<\frac{1}{2}\Rightarrow \frac{\pi}{3}<\Theta<\frac{5\pi}{3}
\]

\[
\Pr\left(\cos\Theta<\frac{1}{2}\right)=\frac{\frac{5\pi}{3}-\frac{\pi}{3}}{2\pi}
=\frac{\frac{4\pi}{3}}{2\pi}=\frac{2}{3}
\]

\hrule
\vspace{1em}

\subsection*{\textbf{2.3 Uniform Random Variable: Application Examples}}

\begin{itemize}
    \item The phase of a sinusoidal signal, when all phase angles between $(0)$ and $(2\pi)$ are equally likely.
    \item A random number generated by a computer between $(0)$ and $(1)$ for simulations.
    \item The arrival time of a user within a known time window, assuming no time preference.
\end{itemize}

\hrule
\vspace{1em}

\section*{\textbf{3. Exponential Random Variable}}

The exponential random variable has a PDF and CDF given by (for any $(b>0)$):

\subsubsection*{\textbf{PDF}}

\[
f_X(x)=\frac{1}{b}\exp\left(-\frac{x}{b}\right)u(x)
\]

\subsubsection*{\textbf{CDF}}

\[
F_X(x)=\left[1-\exp\left(-\frac{x}{b}\right)\right]u(x)
\]

\hrule
\vspace{1em}

\subsection*{\textbf{3.1 Example \#2 (Exponential RV)}}

\subsubsection*{\textbf{Problem}}

Let $(X)$ be an exponential random variable with PDF:

\[
f_X(x)=e^{-x}u(x)
\]

Find:

\begin{itemize}
    \item (a) $(\Pr(3X<5))$
    \item (b) Generalize to find $(\Pr(3X<y))$ for arbitrary constant $(y)$
    \item (c) Let $(Y=3X)$. Find $(f_Y(y))$
\end{itemize}

\hrule
\vspace{1em}

\subsubsection*{\textbf{Solution: Example 2 (Exponential RV)}}

Given:

\[
f_X(x)=e^{-x}u(x)
\]

\[
F_X(x)=[1-e^{-x}]u(x)
\]

\subsubsection*{\textbf{(a) Find $(\Pr(3X<5))$}}

\[
\Pr(3X<5)=\Pr\left(X<\frac{5}{3}\right)
\]

\[
=F_X\left(\frac{5}{3}\right)=1-e^{-5/3}
\]

\subsubsection*{\textbf{(b) Find $(\Pr(3X<y))$}}

\[
\Pr(3X<y)=\Pr\left(X<\frac{y}{3}\right)
\]

\[
=F_X\left(\frac{y}{3}\right)=[1-e^{-y/3}]u(y)
\]

\subsubsection*{\textbf{(c) Let $(Y=3X)$. Find $(f_Y(y))$}}

From part (b), the CDF of $(Y)$ is:

\[
F_Y(y)=[1-e^{-y/3}]u(y)
\]

Differentiating:

\[
f_Y(y)=\frac{d}{dy}F_Y(y)=\frac{1}{3}e^{-y/3}u(y)
\]

Hence, $(Y)$ is an exponential random variable with parameter $\lambda=\frac{1}{3}$.

\hrule
\vspace{1em}

\subsection*{\textbf{3.2 Exponential Random Variable: Application Examples}}

\begin{itemize}
    \item Time until a radioactive particle decays, assuming a constant decay rate.
    \item Duration of an idle period of a wireless communication channel before it becomes busy.
    \item Time until the next earthquake in a simplified seismic activity model with a constant occurrence rate.
\end{itemize}

\hrule
\vspace{1em}

\section*{\textbf{4. Laplace Random Variable}}

\subsection*{\textbf{4.1 Laplace Random Variable}}

\subsubsection*{\textbf{PDF}}

\[
f_X(x)=\frac{1}{2b}\exp\left(-\frac{|x|}{b}\right)
\]

\subsubsection*{\textbf{CDF}}

\[
F_X(x)=
\begin{cases}
\frac{1}{2}\exp\left(\frac{x}{b}\right), & x<0 \\
1-\frac{1}{2}\exp\left(-\frac{x}{b}\right), & x>0
\end{cases}
\]

Application: The Laplace random variable has been used to model the amplitude of a speech (voice) signal.

\hrule
\vspace{1em}

\subsection*{\textbf{4.2 Example \#3 (Laplace RV)}}

\subsubsection*{\textbf{Problem}}

Let $(W)$ be a Laplace random variable with PDF:

\[
f_W(w)=ce^{-2|w|}
\]

Find:

\begin{itemize}
    \item (a) The value of the constant $(c)$
    \item (b) $(\Pr(-1<W<2))$
    \item (c) $(\Pr(W>0 \mid -1<W<2))$
\end{itemize}

\hrule
\vspace{1em}

\subsubsection*{\textbf{Solution: Example 3 (Laplace RV)}}

Given:

\[
f_W(w)=ce^{-2|w|}
\]

\subsubsection*{\textbf{(a) Find $(c)$}}

\[
\int_{-\infty}^{\infty}ce^{-2|w|}dw
=2c\int_0^\infty e^{-2w}dw
\]

\[
=2c\cdot\frac{1}{2}=c
\]

Thus:

\[
c=1
\]

\subsubsection*{\textbf{(b) Find $(\Pr(-1<W<2))$}}

(Only written as given in material.)

\subsubsection*{\textbf{(c) Find $(\Pr(W>0 \mid -1<W<2))$}}

\[
\Pr(W>0\mid -1<W<2)=\frac{\Pr(0<W<2)}{\Pr(-1<W<2)}
\]

\[
\Pr(0<W<2)=\int_0^2 e^{-2w}dw=\frac{1}{2}(1-e^{-4})
\]

\[
\Pr(-1<W<2)=1-\frac{1}{2}(e^{-2}+e^{-4})
\]

\[
\Pr(W>0\mid -1<W<2)=\frac{1-e^{-4}}{2-(e^{-2}+e^{-4})}
\]

\hrule
\vspace{1em}

\subsection*{\textbf{4.3 Laplace Random Variable: Application Examples}}

\begin{itemize}
    \item Modeling impulsive noise in communication systems.
    \item Error distribution in $(L_1)$ regression.
    \item Noise added in differential privacy mechanisms.
\end{itemize}

\hrule
\vspace{1em}

\section*{\textbf{5. Gamma Random Variable}}

\subsection*{\textbf{5.1 Gamma Random Variable}}

\subsubsection*{\textbf{PDF}}

\[
f_X(x)=\frac{(x/b)^{c-1}\exp(-x/b)}{b\Gamma(c)}u(x)
\]

\subsubsection*{\textbf{CDF}}

\[
F_X(x)=\frac{\gamma(c,x/b)}{\Gamma(c)}u(x)
\]

\hrule
\vspace{1em}

\subsubsection*{\textbf{Definitions}}

\paragraph{\textbf{Gamma Function}}

\[
\Gamma(a)=\int_0^\infty e^{-t}t^{a-1}dt
\]

\paragraph{\textbf{Incomplete Gamma Function}}

\[
\gamma(a,x)=\int_0^x e^{-t}t^{a-1}dt
\]

\hrule
\vspace{1em}

\subsection*{\textbf{5.2 Gamma Random Variable (Graphs \& Special Cases)}}

Resulting random variables:

\begin{itemize}
    \item For $(b=2, c=0.5)$, the resulting RV is a Chi-square RV.
    \item For $(c=1)$, the resulting RV is an Exponential RV.
\end{itemize}

where $(k=c)$ and $(\theta=0b)$.

\hrule
\vspace{1em}

\subsection*{\textbf{5.3 Example \#4 (Gamma RV)}}

\subsubsection*{\textbf{Problem}}

Let $(X)$ be a gamma random variable with parameters $(a)$ and $(\lambda)$. Calculate:

\begin{itemize}
    \item (a) $(E[X])$
    \item (b) $(\text{Var}(X))$
\end{itemize}

\hrule
\vspace{1em}

\subsubsection*{\textbf{Solution: Example \#4 (Gamma RV)}}

Given:

\[
X\sim Gamma(a,\lambda)
\]

with PDF:

\[
f_X(x)=\frac{\lambda^a}{\Gamma(a)}x^{a-1}e^{-\lambda x},\quad x>0
\]

\subsubsection*{\textbf{(a) Find $(E[X])$}}

\[
E[X]=\int_0^\infty xf_X(x)dx
\]

\[
=\frac{\lambda^a}{\Gamma(a)}\int_0^\infty x^a e^{-\lambda x}dx
\]

\[
=\frac{1}{\lambda\Gamma(a)}\int_0^\infty e^{-\lambda x}(\lambda x)^a dx
\]

\[
=\frac{1}{\lambda\Gamma(a)}\Gamma(a+1)
\]

\[
=\frac{a}{\lambda}
\]

\subsubsection*{\textbf{(b) Find $(\text{Var}(X))$}}

\[
Var(X)=E[X^2]-(E[X])^2
\]

First compute $(E[X^2])$:

\[
E[X^2]=\frac{\lambda^a}{\Gamma(a)}\int_0^\infty x^{a+1}e^{-\lambda x}dx
\]

\[
=\frac{1}{\lambda^2\Gamma(a)}\int_0^\infty e^{-\lambda x}(\lambda x)^{a+1}dx
\]

\[
=\frac{\Gamma(a+2)}{\lambda^2\Gamma(a)}
\]

\[
=\frac{a(a+1)}{\lambda^2}
\]

Thus:

\[
Var(X)=\frac{a(a+1)}{\lambda^2}-\left(\frac{a}{\lambda}\right)^2
=\frac{a}{\lambda^2}
\]

\hrule
\vspace{1em}

\subsection*{\textbf{5.4 Example \#5 (Gamma RV)}}

\subsubsection*{\textbf{Problem}}

Prove the following properties of the Gamma function:

\begin{itemize}
    \item (a) $(\Gamma(n)=(n-1)!)$ for $(n=1,2,3,...)$
    \item (b) $(\Gamma(x+1)=x\Gamma(x))$
    \item (c) $(\Gamma(1/2)=\sqrt{\pi})$
\end{itemize}

Post your answer on Campuswire!

\hrule
\vspace{1em}

\subsection*{\textbf{5.5 Gamma Random Variable: Application Examples}}

\begin{itemize}
    \item Total service time for completing multiple independent tasks.
    \item Received signal energy in detection systems.
    \item Aggregate waiting time until the $(k)$-th event in a Poisson process.
\end{itemize}

\hrule
\vspace{1em}

\section*{\textbf{6. Exercise: Problem Solving 1 (Movie Theater Optimization)}}

\subsection*{\textbf{Problem}}

Suppose that you plan to arrive at a movie theater for a scheduled showtime at a fixed start time $(T)$.

\begin{itemize}
    \item If you arrive $(s)$ minutes early, waiting cost = $(c\cdot s)$
    \item If you arrive $(s)$ minutes late, missing scene cost = $(k\cdot s)$
    \item Travel time is a continuous RV with pdf $(f(t))$
\end{itemize}

Determine the departure time that minimizes expected cost.

\hrule
\vspace{1em}

\subsection*{\textbf{Solution: Problem Solving 1}}

Let $(d)$ denote departure time from home.\\
Let $(T_r)$ be the random travel time with pdf $(f(t))$ and CDF $(F(t))$.

\subsubsection*{\textbf{Arrival time}}

\[
Arrival=d+T_r
\]

\subsubsection*{\textbf{Cost structure}}

\begin{itemize}
    \item Early arrival $((d+T_r<T))$:
\end{itemize}

\[
cost=c(T-d-T_r)
\]

\begin{itemize}
    \item Late arrival $((d+T_r>T))$:
\end{itemize}

\[
cost=k(d+T_r-T)
\]

\subsubsection*{\textbf{Expected cost}}

\[
E[C(d)]=\int_0^{T-d}c(T-d-t)f(t)dt+\int_{T-d}^{\infty}k(t+d-T)f(t)dt
\]

Differentiate $(E[C(d)])$ with respect to $(d)$ and set to zero:

\[
-cF(T-d)+k(1-F(T-d))=0
\]

\subsubsection*{\textbf{Optimality condition}}

\[
F(T-d)=\frac{k}{c+k}
\]

\subsubsection*{\textbf{Optimal departure time}}

\[
d^*=T-F^{-1}\left(\frac{k}{c+k}\right)
\]

Interpretation: Depart such that the probability of arriving early equals $\left(\frac{k}{c+k}\right)$.

\hrule
\vspace{1em}

\section*{\textbf{7. Exercise: Problem Solving 2 (CDF Analysis)}}

\subsection*{\textbf{Problem}}

The CDF of a RV is:

\[
F_X(x)=
\begin{cases}
0, & x<0 \\
x^2, & 0<x<1 \\
1, & x>1
\end{cases}
\]

Find:

\begin{itemize}
    \item Mean $(\mu_X)$
    \item Variance $(\sigma_X^2)$
    \item Skewness $(c_s)$
    \item Kurtosis $(c_k)$
\end{itemize}

\hrule
\vspace{1em}

\subsection*{\textbf{Solution: Problem Solving 2}}

Given CDF:

\[
F_X(x)=
\begin{cases}
0, & x<0 \\
x^2, & 0<x<1 \\
1, & x>1
\end{cases}
\]

\subsubsection*{\textbf{PDF}}

\[
f_X(x)=\frac{d}{dx}F_X(x)=2x,\quad 0<x<1
\]

\hrule
\vspace{1em}

\subsubsection*{\textbf{Mean}}

\[
\mu_X=\int_0^1 x(2x)dx=\frac{2}{3}
\]

\hrule
\vspace{1em}

\subsubsection*{\textbf{Variance}}

\[
\sigma_X^2=\int_0^1 x^2(2x)dx-\mu_X^2
\]

\[
=\frac{1}{2}-\frac{4}{9}=\frac{1}{18}
\]

\hrule
\vspace{1em}

\subsubsection*{\textbf{Skewness}}

\[
c_s=\frac{E[(X-\mu)^3]}{\sigma^3}=-\frac{2}{5}
\]

\hrule
\vspace{1em}

\subsubsection*{\textbf{Kurtosis}}

\[
c_k=\frac{E[(X-\mu)^4]}{\sigma^4}=\frac{12}{5}
\]

\hrule
\vspace{1em}

\subsubsection*{\textbf{Conclusion}}

\begin{itemize}
    \item Distribution is left-skewed
    \item Platykurtic (lighter tails than Gaussian)
\end{itemize}

\hrule
\vspace{1em}

\section*{\textbf{8. Gaussian Modeling: From Noise to Math}}

\subsection*{\textbf{Motivation: Why Probability? Because the World is Noisy}}

\subsubsection*{\textbf{The Reality}}

\begin{itemize}
    \item Sensors don’t give “flat” lines (thermal noise, interference).
    \item Repeated measurements form a “cloud” of uncertainty.
\end{itemize}

\subsubsection*{\textbf{The Simulation}}

We model noise using the Gaussian RV:

\[
X\sim N(\mu,\sigma^2)
\]

Generative Formula:

\[
X=\sigma Z+\mu
\]

Coding Task 1: Transform standard randomness $(Z)$ into a physical model $(X)$.

Probability models the shape of uncertainty, not just exact values.

\hrule
\vspace{1em}

\section*{\textbf{9. Density Estimation: Learning from Data}}

\subsection*{\textbf{Gaussian Density Estimation from Noisy Observations}}

Reality Check: We don’t know $(\mu)$ or $(\sigma)$ in the field.

\subsubsection*{\textbf{The Problem}}

We only observe raw noisy samples (a histogram).

\subsubsection*{\textbf{The Solution}}

We use sample statistics to estimate the distribution.

Key Concept: We estimate probability distributions, not deterministic functions.\\
This “smooths” the noise into a predictive tool.

\hrule
\vspace{1em}

\section*{\textbf{10. Applications: Network Systems}}

\subsection*{\textbf{Packet Delay Modeling}}

Case Study: 500 ICMP Pings (Latency Analysis)

Scenario:\\
A server transmits packets. Delay varies due to queuing, congestion, and routing hops.

Engineering Goal:

\begin{itemize}
    \item Typical delay $(\mu)$
    \item Jitter (variation) $(\sigma)$
\end{itemize}

Critical Discussion: Why is high jitter $(\sigma)$ more damaging to real-time video than a high constant delay $(\mu)$?

\hrule
\vspace{1em}

\section*{\textbf{11. Applications: Broader Computing Contexts}}

\subsection*{\textbf{Probability in Modern Engineering}}

\subsubsection*{\textbf{Cloud Systems / IoT / Robotics}}

\begin{itemize}
    \item \textbf{Sensor Fusion (GPS noise):} Kalman Filters estimate real location between noisy pings.
    \item \textbf{Image Processing (Denoising):} Gaussian priors separate real signal from thermal noise.
    \item \textbf{Tail Latency (Distributed systems):} model 99th percentile delay for reliability.
\end{itemize}

From hardware sensors to software backends, uncertainty is the only constant.

\hrule
\vspace{1em}

\section*{\textbf{End of Lecture Scribe 9}}

\end{document}
